% !TeX root = ../navier-stokes-manuscript.tex
% ============================================================
% 1.7 Singularity and the Derivative Tower
% ============================================================
\subsection{Singularity and the Derivative Tower}\label{sub:SingularityAndTheDerivativeTower}

The Euler sheet loses smoothness in the velocity itself.
That is the easiest case to see, because the two sides of the interface already disagree on the value of \(u\).
A singularity can occur higher up.
The velocity may stay finite and continuous while its gradient diverges, while two approaches give different accelerations, or while some still higher derivative stops existing.
The first failed quantity can sit anywhere above the visible motion.

The natural place to begin climbing is the acceleration of a material parcel, because it is already written into the Navier--Stokes equation.
Set
\[
  D_t:=\partial_t+u\cdot\nabla.
\]
Along the trajectory \(X(a,t)\), the parcel acceleration is
\[
  a(t,X(a,t))
  :=
  D_tu(t,X(a,t))
  =
  -\nabla p(t,X(a,t))
  +\nu\Delta u(t,X(a,t)).
\]
Pressure and viscosity determine how the parcel velocity changes while transport moves the parcel through the field.

Following that same parcel once more gives its material jerk,
\[
  j:=D_ta=D_t^2u.
\]
To compare this physical derivative with derivatives taken at a fixed point, write
\[
  V_n:=\partial_t^n u.
\]
A direct expansion gives
\[
  D_t^2u
  =
  V_2
  +2(V_0\cdot\nabla)V_1
  +(V_1\cdot\nabla)V_0
  +(V_0\cdot\nabla)\bigl((V_0\cdot\nabla)V_0\bigr).
\]
The fixed-position second time derivative \(V_2\) is one part of the material jerk.
The remaining terms appear because the parcel moves through spatial variation while its velocity is changing.

This distinction stays with us throughout the tower.
Material derivatives \(D_t^ru\) follow a particle path.
Eulerian derivatives \(\partial_t^nu\) remain at one spatial coordinate.
They describe the same evolving field through \(D_t=\partial_t+u\cdot\nabla\), so every material rung expands into Eulerian time derivatives and spatial derivatives of the lower rungs.
The material sequence carries the physical picture, while the Eulerian sequence makes repeated differentiation of the equation transparent.

\divider{The Temporal Rungs}

Write
\[
  P_n:=\partial_t^n p.
\]
The equation itself gives the first Eulerian rung:
\[
  V_1
  =
  \nu\Delta V_0
  -(V_0\cdot\nabla)V_0
  -\nabla P_0.
\]
Differentiate once in time.
The derivative can land on either velocity factor in the transport term, which gives
\[
  V_2
  =
  \nu\Delta V_1
  -(V_1\cdot\nabla)V_0
  -(V_0\cdot\nabla)V_1
  -\nabla P_1.
\]
The next differentiation gives
\[
  V_3
  =
  \nu\Delta V_2
  -(V_2\cdot\nabla)V_0
  -2(V_1\cdot\nabla)V_1
  -(V_0\cdot\nabla)V_2
  -\nabla P_2.
\]
One more gives
\[
  \begin{aligned}
  V_4
  ={}&
  \nu\Delta V_3
  -(V_3\cdot\nabla)V_0
  -3(V_2\cdot\nabla)V_1\\
  &-3(V_1\cdot\nabla)V_2
  -(V_0\cdot\nabla)V_3
  -\nabla P_3.
  \end{aligned}
\]
The coefficients \(1,1\), then \(1,2,1\), then \(1,3,3,1\), appear because the product rule counts every way the time derivatives can split between the two copies of velocity.
At arbitrary order,
\[
  V_{n+1}
  +
  \sum_{k=0}^{n}
  \binom{n}{k}
  (V_k\cdot\nabla)V_{n-k}
  +
  \nabla P_n
  =
  \nu\Delta V_n,
  \qquad
  \nabla\cdot V_n=0.
\]
Every temporal rung is connected to the lower rungs through transport, to spatial curvature through viscosity, and to pressure at the same time order.
The higher responses have not been added from outside the equation.
They are generated by differentiating the equation the fluid already obeys.

\divider{Pressure and the Spatial Rungs}

Pressure follows the velocity tower at every order.
Taking the divergence of the original momentum equation and using \(\nabla\cdot u=0\) gives
\[
  -\Delta p
  =
  \partial_i\partial_j(u_i u_j),
\]
where repeated spatial indices are summed.
After \(n\) time derivatives,
\[
  -\Delta P_n
  =
  \partial_i\partial_j
  \sum_{k=0}^{n}
  \binom{n}{k}
  V_{k,i}V_{n-k,j}.
\]
The velocity rungs determine \(P_n\) through an elliptic problem over \(\mathbb R^3\).
The pressure response at one location can depend on the velocity configuration far beyond an arbitrarily small neighbourhood of that location.

Spatial differentiation opens the other direction of the same structure.
For one coordinate direction \(x_\ell\),
\[
  \begin{aligned}
  \partial_t(\partial_\ell u)
  &+
  (u\cdot\nabla)(\partial_\ell u)
  +
  ((\partial_\ell u)\cdot\nabla)u
  +
  \nabla(\partial_\ell p)\\
  &=
  \nu\Delta(\partial_\ell u).
  \end{aligned}
\]
At the next spatial rung the product rule produces both cross terms:
\[
  \begin{aligned}
  \partial_t(\partial_{\ell m}u)
  &+
  (u\cdot\nabla)(\partial_{\ell m}u)
  +
  ((\partial_\ell u)\cdot\nabla)(\partial_m u)\\
  &+
  ((\partial_m u)\cdot\nabla)(\partial_\ell u)
  +
  ((\partial_{\ell m}u)\cdot\nabla)u
  +
  \nabla(\partial_{\ell m}p)\\
  &=
  \nu\Delta(\partial_{\ell m}u).
  \end{aligned}
\]
Every further time or spatial derivative repeats this same nesting.

Using a multi-index \(\alpha\) for the spatial derivatives, collect all of the velocity values at \((t,x)\) in
\[
  \mathcal J^\infty u(t,x)
  :=
  \left\{
    \partial_t^n\partial_x^\alpha u(t,x):
    n\geq0,\ 
    \alpha\in\mathbb N_0^3
  \right\}.
\]
This is the derivative tower.
Its base is the velocity, its vertical direction contains the fixed-position time derivatives, and its spatial directions contain the gradients and curvatures that enter transport and viscosity.
The pressure tower remains attached through the differentiated Poisson equations.
Material acceleration, jerk, and every higher material response are combinations drawn from this mixed structure along a moving trajectory.

Individual rungs are allowed to vanish.
The zero solution has a vanishing tower, and a steady flow has \(V_n=0\) for every \(n\geq1\), while both can remain smooth.
What smoothness requires is that every finite rung under consideration remain finite, approach one compatible value, and continue to be a derivative of the same classical field.
A proposed endpoint can lose one of those properties at the velocity, acceleration, jerk, a spatial derivative, or any mixed rung above them.
